\part{Computation as Rewrite}
\label{part:rewrite}

\chapter*{Introduction to Part I}
\addcontentsline{toc}{chapter}{Introduction to Part I}

\begin{quote}
\emph{Part I builds the foundation. At this layer, we stay disciplined.  
We are talking about engineered systems\emdash{}software, pipelines, state, structure, motion.  
Nothing more exotic is assumed.}
\end{quote}

\bigskip

\section*{Why Rewrite?}
\label{sec:why-rewrite}

The Prologue showed a recurring pattern in modern systems:  
everything we build eventually collapses into structure and change.

Part I formalizes that pattern.

To work honestly at this layer, we need a lens that matches the world we actually engineer.  
Not the comforting metaphors of the 1970s\emdash{}files, threads, stacks\emdash{}but the deeper
geometry underneath them: \textbf{state evolving under rules}.

In this part of the book, we take a disciplined position:

\begin{quote}
\textbf{Computation is best understood as rewrite:  
the transformation of structured state under explicit rules.}
\end{quote}

That’s it.  
No cosmology. No metaphysics. No claims about “the universe.”  
Those questions will come later, once we have earned the right to ask them.

For now, we are building the formal substrate for everything that follows.

\bigskip

\section*{A Working Hypothesis}
\label{sec:working-hypothesis}

The systems you work on today\emdash{}distributed backends,
data pipelines, build graphs, reactive UIs, real-time engines\emdash{}do not behave like lists of instructions.
They behave like \textbf{structured state undergoing transformation}.

Rewrite is not a metaphor.  
Rewrite is the geometry of the motion.

This part of the book establishes three ideas:

\begin{enumerate}
    \item \textbf{Structure:} why engineered systems naturally collapse into graph form.
    \item \textbf{Rewrite:} how state evolves under rules.
    \item \textbf{Composition:} how these rewrites build, layer, and interact.
\end{enumerate}

These ideas are the intellectual spine of the entire book.

\bigskip

\section*{What This Part Is Not}
\label{sec:not-yet}

You may notice echoes\emdash{}between rewrite and physics,
rewrite and cognition, rewrite and dynamical systems.
That is intentional.

But we are not making claims about the universe.

Not yet.

In Parts IV–VI, we will explore how far this geometry reaches.
For now, we stay grounded: the domain is engineered systems.

Part I is about \textbf{what software already is},  
not what the universe might be.

\bigskip

\section*{The Road Ahead}
\label{sec:road-ahead}

Part I sits at the base of a long ascent:

\begin{itemize}
    \item Chapters 2–4 build the substrate: Recursive Meta-Graphs, rewrite rules,
    causality, compositional structure.
    \item Part II explores movement: execution as path, history as geometry.
    \item Part III introduces distance, curvature, and possibility.
    \item Part IV builds machines atop that geometry.
    \item Part V connects it all to correctness, concurrency, and provenance.
    \item Part VI finally asks the forbidden question:  
    \begin{center}
        \emph{If rewrite explains engineered systems so elegantly\emdash{}  
        what else might it describe?}
    \end{center}
\end{itemize}

But all of that begins here, with a single disciplined idea:

\[
\text{state} \xrightarrow{\text{rewrite}} \text{state}'.
\]

This is computation’s most primitive motion.  
Everything else is structure on top of it.

    